\section{Executable polynomial presentation}

For $A\in M_n(\Q)$ we define
\[
  \operatorname{presentation}(A)_{ij} =
  \begin{cases}
    X-A_{ii}, & i=j,\\
    -A_{ij}, & i\ne j.
  \end{cases}
\]
The Lean definition \code{endomorphismPresentationMatrix} is executable and a
separate theorem identifies it exactly with
\code{Matrix.charmatrix}.  This bridge fixes the sign convention once and
allows later statements to use mathlib's characteristic polynomial API.

The core polynomial algorithm runs over an explicit finite-support runtime
whose rational arithmetic, division, normalization, and extended gcd are
constructive.  A coefficient-preserving ring equivalence maps runtime
polynomials to standard $\Q[X]$.  Mapping a strong Smith certificate across
this equivalence gives
\code{endomorphismSmithResult}: it retains $U$, $U^{-1}$, $S$, $V$, $V^{-1}$,
the exact transformation equation, and the Smith predicate.

For each $i:\operatorname{Fin} n$ we read
\[
  d_i = S_{ii}.
\]
The public list \code{endomorphismInvariantFactors} contains all $n$ entries.
Keeping units is important: it preserves the exact square Smith diagonal and
makes list length, product, and dependent indexing statements uniform.
Human-facing output may filter units later, but theorem data does not.

The Smith predicate proves that every $d_i$ is normalized and that consecutive
entries divide.  Over $\Q[X]$, normalized nonzero factors are monic.  A
determinant calculation using the unit determinants of $U$ and $V$ first gives
association between $\det S$ and $\det(XI-A)$.  Normalization on both sides
strengthens this to the exact identity
\[
  \prod_i d_i = \chi_A.
\]
Dimension zero is not patched by an auxiliary assumption: the factor list is
empty, its product is one, and the option-valued largest-factor reader returns
\code{none}.
